Выполненная выпускная квалификационная работа посвящена разработке генератора
данных для аналитического бенчмарка TPC-H, ориентированного на современную
аналитическую инфраструктуру. В ходе работы проанализирована предметная
область, рассмотрены существующие генераторы, сформулированы требования к
масштабируемой архитектуре и реализована программная система, поддерживающая
детерминированное формирование таблиц, подготовку колонoчных батчей и запись
результата в современные форматы хранения.

Поставленные задачи решены в полном объеме. Разработана архитектура,
обеспечивающая параллельную генерацию на основе быстрой перемотки линейного
генератора случайных чисел, реализован слой сериализации для нескольких
колонoчных форматов, обеспечена работа с локальными и удаленными хранилищами,
а также показана возможность использования генератора без файлового слоя в
качестве встроенного источника данных для аналитической СУБД.

Практическая значимость результата состоит в сокращении числа промежуточных
этапов подготовки набора данных. Генератор позволяет сразу получать данные в
колонoчном представлении и целевых форматах хранения, что упрощает
воспроизводимое проведение экспериментов и снижает накладные расходы на
подготовку данных. Работоспособность встроенного сценария подтверждена
реализацией отдельного расширения для DuckDB, выступающего как
\textit{proof-of-concept} использования генератора без слоя сериализации.

С точки зрения научно-технического уровня разработанное решение сочетает
детерминированность эталонного TPC-H с прямой генерацией колонoчных батчей,
поддержкой современных форматов хранения и возможностью встроенного
использования. Это позволяет рассматривать полученный результат как основу для
дальнейшего развития генератора в направлении новых форматов и новых
сценариев интеграции с аналитическими системами.
